\section{Related Work}
\label{sec:rel}

The literature review mainly focuses on next-best-view methods that are formulated in 5 DoF or assume a drone as an agent.\\
\textbf{Scene-specific next-best-view planners.} Next-best-view planners have demonstrated promising results in active 3D reconstruction by predicting the optimal viewpoint for data capture based on the current state.
%(e.g., the reconstruction progress).
Traditional approaches~\cite{monica2018contour, zha1997next, liu2018object, hardouin2020surface, hardouin2020next, dai2020fast, zhou2021fuel} rely on hand-crafted rules to determine the next-best viewpoint. 
For instance, the method~\cite{zha1997next} selected the next-best viewpoint by maximizing a rating function favoring smooth regions, which may overlook fine-grained object details.
Instead, the methods~\cite{monica2018contour, hardouin2020surface} collected data along boundaries between seen and unseen surfaces to capture finer details, but still require handcrafted parameter tuning for each scene.
Another approach~\cite{liu2018object} scanned segmented objects sequentially using a predefined object database, reducing the need for handcrafted tuning, but its performance degrades in cluttered environments due to inaccurate object matching.
Recent advances in deep learning and increased computational power have given rise to learning-based methods~\cite{lin2022active, pan2022activenerf, lee2022uncertainty, zhan2022activermap, jin2023neu, sunderhauf2023density, chen2024gennbv, peralta2020next, ran2023neurar, zeng2022efficient, tao2022seer, wilson2025pop, li2025activesplat}.
Some studies~\cite{lin2022active, sunderhauf2023density} used NeRF~\cite{mildenhall2020nerf} ensembles to estimate uncertainty via model disagreement for viewpoint selection, resulting in linearly increasing computational overhead.
Other works~\cite{pan2022activenerf, ran2023neurar, zeng2022efficient} avoid the computational overhead by incorporating Bayesian-based NeRF~\cite{shen2021stochastic, martin2021nerf} to estimate uncertainty for viewpoint selection.
Meanwhile, other approaches~\cite{lee2022uncertainty, zhan2022activermap} defined the next-best viewpoint as the viewpoint that maximizes the entropy of the density field along the camera rays.
Although these methods have shown outstanding performance in collecting data, they typically require sampling candidate viewpoints.
In addition, their reliance on online learning makes them less suitable for real-time applications. \\
\textbf{Generalizable next-best-view planners.} 
Unlike the aforementioned online-learning approaches, the generalizable methods~\cite{jin2023neu, chen2024gennbv, peralta2020next} avoid the training process for new scenes, thereby enabling faster next-best-view selection. 
Prediction time is important for real-world tasks where a robot may run out of battery within a few minutes.
Among generalizable methods, prior work~\cite{jin2023neu} proposed a Bayesian-based NeRF that selects the next-best viewpoint to maximize view variance without additional training.
However, it still requires candidate viewpoint sampling, resulting in performance unsuitable for real-time applications.
Instead, another line of generalizable next-best-view approaches~\cite{chen2024gennbv, peralta2020next} utilized reinforcement learning to learn a next-best-view planner to bypass the need for sampling candidate viewpoints.
Prior work~\cite{peralta2020next} proposed learning a 3-DoF next-best-view policy using a series of grayscale images as observations.
Subsequently, prior work~\cite{chen2024gennbv} improved upon this method by incorporating occupancy grids into the observations, which provide explicit geometric information. 
This enhancement enabled the development of a 5-DoF next-best-view planner, achieving an outstanding coverage ratio for unknown scenes.

As shown in Tab.~S4, compared to the scene-specific methods, Hestia does not require candidate viewpoint sampling or inference-time optimization, thereby enabling more flexible viewpoint prediction. 
Compared to the generalizable method~\cite{chen2024gennbv}, Hestia treats voxels as cubes rather than points, enabling more comprehensive capture. 
In addition, Hestia adopts a close-greedy scheme to mitigate spurious correlations, introduces a hierarchical structure to model the action space, and uses a more diverse dataset to maintain robust performance across varying object shapes and positions.
Notably, Hestia’s hierarchical structure addresses the challenge of high-dimensional continuous action search in RL-based generalizable next-best-view planning, which differs from traditional methods~\cite{monica2018contour, hardouin2020surface, dai2020fast, zhou2021fuel}.
%Notably, Hestia’s hierarchical structure addresses the challenge of high-dimensional continuous action search in RL-based generalizable next-best-view planning. 
%This differs from traditional methods~\cite{monica2018contour, hardouin2020surface, dai2020fast, zhou2021fuel}, which use hierarchical structures to traverse frontiers or optimize tour efficiency.